\chapter{Literature and Theoretical Background LP}
\label{ch:literature1}

This chapter reviews the literatures on which the empirical strategy
of the thesis builds and locates the thesis's contribution within
them.

The chapter is organized argumentatively rather than chronologically.
\Cref{sec:literature-hfi} introduces the high-frequency-identification
approach and the empirical finding that a single target-rate surprise
is insufficient to summarize asset-price responses to FOMC announcements,
which motivates the use of a forward-guidance shock as the primary
empirical object. \Cref{sec:literature-forward-guidance} examines the
forward-guidance dimension specifically and the central-bank
information-effect concerns that complicate its interpretation; this
section also motivates the path-information split used as a robustness
check later in the thesis. \Cref{sec:literature-equity-valuation}
turns to the equity-valuation literature, establishing how monetary
surprises move equity prices through the discount-rate, cash-flow,
and risk-premium channels of a present-value decomposition, and how
firms with growth-firm characteristics are differentially exposed to
the discount-rate component. \Cref{sec:literature-firm-heterogeneity}
reviews the firm-level heterogeneity literature on monetary
transmission, which establishes that responses vary systematically
along financial and structural firm characteristics without yet
incorporating AI exposure as one of those characteristics.
\Cref{sec:literature-ai-intangibles} develops the asset-pricing
bridge from the AI-investment literature to the duration-based view
of equity valuation, and is the section in which the thesis's
mechanism narrative is assembled. \Cref{sec:literature-gap} states
the gap that the thesis fills and the boundaries of the contribution.

Three framing points apply throughout. First, the thesis's measure of
AI exposure is textual, constructed from pre-2023 SEC 10-K filings;
it captures disclosed AI engagement rather than AI adoption directly,
and the literature reviewed below is used to motivate the design
rather than to validate any specific equivalence between text-based
exposure and underlying technology use. Second, the empirical design
estimates differential equity-return responses by AI exposure group
relative to a No AI Exposure reference group, within sector and FOMC
event; the literature on average equity responses to monetary
surprises is relevant for sign and magnitude intuition, but the
thesis's identification operates at the differential level rather
than the level of average effects. Third, no single paper in the
reviewed literature documents the full chain from AI exposure to
differential forward-guidance sensitivity. The mechanism narrative
developed in \cref{sec:literature-ai-intangibles} is a synthesis
across four literatures rather than a single causal claim, and the
empirical results in subsequent chapters should be read accordingly.

\section{Monetary-policy shocks and high-frequency identification}
\label{sec:literature-hfi}

Measuring the causal effect of monetary policy on asset prices requires
isolating the unexpected component of policy announcements from
information that markets had already priced in. The high-frequency
identification (HFI) approach addresses this problem by measuring
changes in interest-rate-sensitive financial instruments inside narrow
windows around scheduled Federal Open Market Committee (FOMC)
announcements. The identifying assumption is that, within a sufficiently
short window, the only material news affecting these instruments is the
content of the FOMC announcement itself, so that the in-window price
change identifies the surprise component of the policy decision
\citep{gurkaynakSackSwanson2005}. This logic is consistent with the
efficient-markets premise that public information is rapidly impounded
into prices \citep{fama1970}, and it has become the dominant empirical
strategy for measuring monetary-policy effects on financial markets.

A central empirical finding of this literature is that a single
target-rate surprise is insufficient to summarize the asset-price
response to FOMC announcements. \citet{gurkaynakSackSwanson2005} show
that two factors are needed to characterize the joint movement of
short- and longer-dated interest-rate futures around FOMC events: one
factor captures the contemporaneous surprise to the federal funds
target, and a second, orthogonal factor captures news about the
expected future path of the policy rate that is conveyed by FOMC
statements rather than by the immediate rate decision. The second
factor accounts for a substantial share of the variation in
longer-dated yields around announcement windows, indicating that
``statements'' --- not only ``actions'' --- carry quantitatively
important market-moving information.

\citet{swanson2021} extends this framework to the post-2008 environment
in which the federal funds rate was constrained at the effective lower
bound for extended periods. Using FOMC announcements from 1991 through
2019, he identifies three orthogonal factors --- a federal funds rate
surprise, a forward-guidance surprise, and a large-scale asset purchase
(LSAP) surprise --- and documents that each factor has economically
large and statistically significant effects on yields across the term
structure, on equity indices, and on exchange rates. The forward-guidance
factor in \citet{swanson2021} is the empirical object that this thesis
uses as its primary shock, and its construction directly motivates the
use of a path-shock series rather than a target-rate series for studying
medium-horizon equity responses. \Cref{sec:literature-forward-guidance}
returns to the interpretation of forward-guidance shocks and to the
information-effect concerns that complicate that interpretation.

A separate but related line of work establishes that high-frequency
surprises carry macroeconomic content and not only financial-market
noise. \citet{gertlerKaradi2015} use high-frequency FOMC surprises as
external instruments in a structural vector autoregression and find
that the identified shocks generate output, price-level, and credit-spread
responses consistent with standard monetary transmission. This evidence
matters for the present thesis because it supports interpreting movements
in HFI-identified shocks as economically meaningful policy news, rather
than as residual variation orthogonal to the macroeconomy.

% [OPTIONAL — keep or cut. Brief signpost only; full LP discussion in Chapter 4.]
While the empirical literature on monetary-policy effects on asset
prices originated in event-window settings, recent work increasingly
estimates dynamic responses horizon by horizon using local projections
\citep{jorda2005,jordaTaylor2025}. \Cref{ch:empirical-strategy} discusses
the local-projection design used in this thesis; for the purposes of the
literature review it is sufficient to note that the identification of
the shock and the estimation of its dynamic effects are separable steps.

Two qualifications to this literature are important for what follows.
First, the validity of high-frequency identification depends on the
window being narrow enough that no other material news arrives within
it, an assumption that is more plausible for scheduled FOMC meetings
than for unscheduled communications \citep{gurkaynakSackSwanson2005}.
The empirical sample in this thesis is accordingly restricted to
scheduled FOMC events. Second, even within a clean announcement
window, the identified surprise need not reflect a pure shift in the
stance of monetary policy: it may also reflect the central bank's
revelation of private information about the economic outlook, or the
market's updating of its beliefs about the policy rule. These concerns
are addressed in \cref{sec:literature-forward-guidance}.

\section{Forward guidance, path shocks, and information-effect concerns}
\label{sec:literature-forward-guidance}

The forward-guidance dimension of monetary-policy surprises is central
to this thesis because it is the empirical object that the path shock
\citep{swanson2021} is designed to measure: news about the expected
future path of the policy rate that is conveyed by FOMC statements,
beyond the contemporaneous surprise to the federal funds target. This
dimension is empirically large --- statement-based factors account for
a substantial share of the market response at FOMC announcements
\citep{gurkaynakSackSwanson2005,swanson2021} --- and it is also the
dimension most vulnerable to the identification concern that high-frequency
surprises do not always reflect pure monetary-policy news. This section
reviews that concern and the empirical strategies that have been proposed
to address it.

The central challenge is that an FOMC announcement bundles a policy
decision with new information conveyed by the central bank about the
economic outlook. \citet{nakamuraSteinsson2018} document that
high-frequency surprises within FOMC windows move not only nominal
yields but also real interest rates and survey-based forecasts of
output and inflation. They interpret this co-movement as evidence that
markets infer information about the state of the economy from FOMC
announcements --- a ``Fed information effect'' --- which implies that
the same measured surprise can reflect a mix of policy news and
disclosure of the central bank's private outlook. For forward-guidance
shocks specifically, this concern is sharper than for contemporaneous
target-rate shocks, because forward guidance is conveyed primarily
through the statement language that markets read as informative about
both the policy rule and the macroeconomic environment
\citep{cieslakSchrimpf2019,hansenMcMahonTong2019}.

\citet{jarocinskiKaradi2020} propose an identification strategy that
separates monetary-policy shocks from information shocks using the
sign of the joint movement of interest rates and broad equity indices
inside the announcement window. The intuition is that a contractionary
monetary-policy shock raises rates and lowers equity prices, whereas a
positive disclosure about the economic outlook raises both rates and
equities. Sign restrictions on this co-movement therefore identify
two distinct shocks within the announcement window. Applying this
decomposition to FOMC and ECB announcements, they find that a
non-trivial share of the variation in conventional high-frequency
surprises reflects information shocks rather than policy shocks, and
that the macroeconomic responses of the two components differ in
direction and persistence. The Jarociński--Karadi sign restriction
provides the empirical template for the path-information split that
this thesis uses as a robustness check.

\citet{bauerSwanson2023} offer an alternative interpretation of the
empirical evidence that has been read as a Fed information effect,
which they label the ``Fed response to news'' channel. They show
that economic news released between the timing of professional
forecast surveys and the subsequent FOMC announcement is correlated
with both the forecast revision and the high-frequency monetary
policy surprise, so that standard regressions used to detect a Fed
information effect suffer from omitted-variables bias; once such
news is controlled for, the puzzling positive coefficient reverses
sign. They rationalise this pattern with a model in which the
private sector has incomplete knowledge of the Fed's monetary policy
rule and therefore underestimates ex ante how strongly the Fed will
respond to publicly available macroeconomic news. Under this
interpretation, the co-movement that \citet{nakamuraSteinsson2018}
and \citet{jarocinskiKaradi2020} attribute to information shocks
reflects a different economic object --- the joint response of the
Fed and private forecasters to common public news, combined with
revisions to the perceived policy rule --- rather than disclosure of
the central bank's private outlook. For the present thesis, the
practical implication of this debate is that the path-information
split should be interpreted with appropriate caution: it isolates
the component of forward-guidance surprises that co-moves with
equities in a way consistent with non-policy news, but the
structural label attached to that component is contested.

\citet{jarocinski2024} provides further refinement by decomposing FOMC
surprises in the unconventional-policy period into multiple components,
including a standard monetary-policy shock, an Odyssean forward-guidance
shock (commitments about the future policy path), a Delphic forward-guidance
shock (revelations about the central bank's outlook), and an LSAP shock.
This decomposition reinforces the more general point that a
forward-guidance surprise measured from announcement-window co-movement
is a composite object whose internal mix of policy news and information
news can vary across episodes. The path shock used in this thesis
\citep{swanson2021} does not distinguish these sub-components; instead,
the empirical strategy addresses the concern at the regression stage by
running the path-information split as a robustness check
(\cref{ch:robustness}).

Three implications follow for the empirical design of this thesis.
First, the path shock is the appropriate primary shock for studying
medium-horizon equity responses because it captures the dimension of
monetary news with the largest effect on longer-dated yields and on
discount-rate-sensitive equity prices; this is the conclusion of
\citet{gurkaynakSackSwanson2005} and \citet{swanson2021}, and it
motivates the use of \texttt{path\_3f} as the main regressor.
Second, the same identification properties that make the path shock
informative also make it the surprise dimension most exposed to the
information-effect concern; a credible empirical specification must
therefore acknowledge this caveat and provide a robustness check that
mitigates --- though does not resolve --- the ambiguity. Third, the
appropriate robustness check, following \citet{jarocinskiKaradi2020},
is to split path surprises by the sign of the contemporaneous equity
reaction within the announcement window, retaining those with
opposite-sign rate and equity movements as the more monetary-leaning
subset and treating same-sign movements as the more information-leaning
subset. This is precisely the construction implemented in the thesis,
and the interpretation of the heterogeneity results in
\cref{ch:results} is conditioned accordingly.

\section{Monetary policy and equity valuation}
\label{sec:literature-equity-valuation}

The previous sections motivated the use of a forward-guidance/path
shock and discussed the identification concerns that surround it. This
section turns to the second component of the empirical strategy: the
outcome variable. The thesis estimates the response of firm-level
cumulative log returns to path shocks, and the relevant literature is
the body of work that documents how monetary-policy surprises move
equity prices and through which channels they do so.

\Citet{bernankeKuttner2005} provide the empirical benchmark for the
average response of broad equity indices to FOMC surprises. Using
high-frequency federal funds rate surprises and a Campbell--Ammer
log-linear return decomposition, they document that an unanticipated
25-basis-point cut in the federal funds rate target is associated
with roughly a one percent increase in broad U.S. equity indices on
average. Decomposing this response into news about real interest
rates, future dividends, and expected future excess returns
(risk premia), they find that the dominant channel is news about
expected future excess returns, with revisions to expected dividends
contributing a secondary and sample-dependent share, and revisions
to expected real interest rates playing a negligible direct role.
The \citet{rigobonSack2004} heteroskedasticity-based identification
delivers a quantitatively similar response using a different
identification strategy, lending robustness to the basic empirical
fact that monetary-policy surprises move equity prices in the
expected direction and with economically meaningful magnitudes. The
sign convention adopted in this thesis follows directly from this
literature: a positive path surprise corresponds to a hawkish/tightening
surprise, and a negative coefficient on the path-shock interaction
therefore means the affected firms fall by more than the reference
group.

A separate body of work decomposes equity-return variation into news
about expected cash flows and news about discount rates.
\Citet{campbellAmmer1993} apply a present-value identity to monthly
U.S. equity and bond returns and show that variation in long-horizon
equity returns can be attributed to a mix of expected-cash-flow news,
expected-real-rate news, expected-inflation news, and expected-excess-return
(risk-premium) news. This decomposition is not a monetary-shock event
study; rather, it provides the accounting framework within which the
average equity response to a monetary surprise can be interpreted as
the sum of contributions from these channels. \Citet{campbellVuolteenaho2004}
develop this framework further by distinguishing ``cash-flow beta''
from ``discount-rate beta'': two assets with the same overall market
exposure can have very different sensitivities to discount-rate news
versus cash-flow news, and these two betas command different risk
premia in the cross-section. \Citet{campbellPolkVuolteenaho2010} apply this decomposition to the growth--value cross-section and document that growth stocks carry
disproportionately high ``good betas'' with market discount-rate
news, whereas value stocks carry disproportionately high ``bad
betas'' with market cash-flow news. They further show that these
patterns are driven by the cash-flow fundamentals of growth and
value firms rather than by transitory shifts in investor sentiment,
which sharpens the asset-pricing interpretation: growth firms' high
sensitivity to market discount-rate news is a property of their
cash-flow structure, not of their pricing alone. This finding is
central to the mechanism narrative developed in
\cref{sec:literature-ai-intangibles}, because it provides the
asset-pricing channel through which firms with growth-firm
characteristics --- long-duration cash flows, large intangible
capital, prominent growth options --- are expected to be more
sensitive to news about expected future policy rates.

\Citet{vanBinsbergen2020} formalizes the closely related
duration-based view of equity valuation. In a present-value framework,
the sensitivity of a stock's price to a change in the discount rate is
increasing in the average maturity of its expected cash flows, in
direct analogy with bond duration. Stocks whose value derives
disproportionately from cash flows far in the future are mechanically
more sensitive to changes in the expected discount-rate path than
stocks whose value derives mostly from near-term cash flows. For the
purposes of this thesis, the relevant implication is that path
surprises --- which move expectations about the policy rate over
multiple quarters and years --- are the dimension of monetary news
that should differentially affect long-duration equities, conditional
on the channels through which the path shock is transmitted to
discount rates being effectively present in the data.

A central limitation of the discount-rate-only narrative is that it
attributes the equity response to monetary surprises to a single
channel, whereas the present-value identity in \citet{campbellAmmer1993}
admits at least three. Recent work has reopened the question of which
channel is quantitatively dominant. \Citet{acostaEtAl2025} use dividend
strips, inflation swaps, and the term structure of interest rates to
decompose the equity response to FOMC communication surprises and
find that revisions to expected risk premia and to expected dividends
contribute substantially to the negative equity response, alongside
the standard discount-rate channel. For this thesis, the implication
is that the heterogeneity results --- if found --- should not be
interpreted as evidence of a single mechanism. Differential responses
of AI-exposed firms relative to non-AI firms within sector and event
could reflect differences along any of the channels in the
\citet{campbellAmmer1993} decomposition, and the empirical design does
not separately identify them.

A smaller body of work studies firm-level equity responses to FOMC
announcements directly, and is therefore the closest empirical
parallel to the design of this thesis. \Citet{gurkaynakKarasoyCanLee2022}
estimate firm-level stock-price responses to high-frequency monetary
surprises and show that the cross-sectional variation in these
responses is systematically related to firms' debt structure ---
the type of debt and its maturity --- and to the forward-guidance
component of the surprise. The paper provides two pieces of evidence
that are directly relevant to the present thesis. First, it
establishes that firm-level equity reactions to FOMC surprises are
sufficiently heterogeneous and well-measured to be analyzed
cross-sectionally, which validates the basic feasibility of a
firm-level event-panel design. Second, it documents that
forward-guidance surprises in particular interact with firm
characteristics in shaping the equity response, which is the type
of interaction the present thesis estimates using AI exposure as the
firm characteristic. The literature does not, however, study AI
exposure as a heterogeneity dimension, and the cumulative-return
horizon studied here (h up to sixty trading days) is longer than the
announcement-window response that is the primary focus of
\citet{gurkaynakKarasoyCanLee2022}; both differences are revisited in
\cref{sec:literature-gap}.

\section{Firm-level heterogeneity in monetary-policy transmission}
\label{sec:literature-firm-heterogeneity}

The literature on firm-level monetary transmission documents that the
response of firms to monetary-policy shocks varies systematically with
observable firm characteristics. This is the strand of work in which
the heterogeneity hypothesis of the present thesis is most naturally
situated: AI exposure is proposed as an additional dimension along
which firm-level responses differ, complementing but not replacing the
financial and structural characteristics that the existing literature
has already identified. This section reviews that evidence, with two
emphases. First, the literature is not uniform in the direction of the
heterogeneity it finds, and the apparent tensions are informative for
how the present thesis should interpret its own results. Second, almost
none of this work uses AI exposure as the heterogeneity dimension, and
the small number of papers that study firm-level equity responses to
forward-guidance shocks do so along orthogonal characteristics. The
gap that the present thesis addresses is therefore one of
characteristic, not of method or design.

A widely cited firm-panel finding is that monetary-policy shocks have
larger real effects on firms that look constrained on conventional
measures. \Citet{cloyneEtAl2023} estimate firm-level investment
responses to monetary surprises in U.S. data and find that younger
firms and firms that do not pay dividends respond more strongly than
older or dividend-paying firms, a pattern they interpret as consistent
with a financial-accelerator mechanism in which weaker access to
external finance amplifies the transmission of monetary policy.
\Citet{durante2022} report a complementary pattern in a panel of
roughly nine million European firm-year observations: younger firms
respond more, firms with higher leverage respond more, and firms in
durable-goods sectors respond more, with high-frequency identified
monetary shocks as the source of variation.

A counterweight to this constraints-amplify-transmission narrative
comes from \citet{ottonelloWinberry2020}. They estimate firm-level
investment responses to monetary shocks in U.S. Compustat data and
document the opposite pattern in the dimension of default risk:
firms with low leverage, high distance to default, and high credit
ratings respond more strongly than firms with high default risk on
the same measures. They reconcile this finding with a
heterogeneous-firm New Keynesian model with default risk in which
high-default-risk firms face a steeper marginal cost curve for
financing investment, while low-risk firms face a flat region of
that curve. A monetary shock shifts out the marginal benefit of
investment; high-risk firms, sitting on the steep part of the
marginal cost curve, absorb part of the shock into higher
borrowing costs rather than into higher investment, dampening their
response relative to low-risk firms. The \citet{ottonelloWinberry2020}
result does not contradict the broader financial-heterogeneity
literature so much as it qualifies it: ``more financially
vulnerable'' is not synonymous with ``more responsive to monetary
policy.'' Reading across this literature, the direction of the
heterogeneity appears to depend on which margin a given firm is
operating against --- the investment opportunity that monetary
policy stimulates, the external-finance cost it works through, or
the proximity to default that bounds it above --- and the
AI-exposure heterogeneity studied below is proposed as a further
such dimension.

A related body of work refines the analysis by separating debt
structure from debt levels. \Citet{caglioEtAl2021} use loan-level and
firm-level data to show that highly leveraged firms and firms that
pledge collateral respond more strongly to monetary surprises through
the credit channel. \Citet{deng2022} report that debt maturity is an
independent dimension of heterogeneity: firms with shorter-maturity
debt respond more strongly to monetary surprises, independently of
total leverage, consistent with rollover risk shaping the transmission
of short-rate news to firm-level investment. \Citet{bahajEtAl2022}
study employment rather than investment as the real outcome and find
that younger and more leveraged firms in the United Kingdom respond
more strongly to monetary shocks, with residential collateral providing
a measurable transmission channel. Taken together, these papers
establish that the heterogeneity literature is multi-dimensional ---
leverage, maturity, collateral, age, dividend status, and sector all
enter --- and that the marginal contribution of any new heterogeneity
dimension must be evaluated against this multidimensional benchmark.

A smaller body of work studies firm-level \emph{equity} responses to
monetary surprises, which is the most direct precedent for the design
of this thesis. \Citet{arinEtAl2025} estimate cross-firm variation in
equity responses to identified conventional and unconventional monetary
shocks and find that leverage is the dominant characteristic explaining
heterogeneity, more so than firm size or age. The firm-level
event-study precedent set by \citet{gurkaynakKarasoyCanLee2022},
already introduced in \cref{sec:literature-equity-valuation},
strengthens the same point in the direction of debt structure: the
type of debt and its maturity, interacted with the forward-guidance
component of the surprise, materially shape the equity response. The
present thesis builds on this firm-level equity-response design and
extends it along a characteristic dimension that this literature has
not used.

Two implications follow for the present thesis. First, AI exposure
should not be presented as a substitute for the financial-heterogeneity
dimensions that the existing literature has established; it is at most
a complementary dimension, and any heterogeneity attributed to AI
exposure must be defended against the alternative explanation that it
reflects an omitted financial characteristic correlated with AI
exposure in the cross-section. The thesis addresses this in
\cref{ch:robustness} by re-estimating the preferred specification on
the SEC-debt subsample and by interacting the path shock with debt-to-assets
on top of AI exposure; the specifications are documented in
\cref{ch:empirical-strategy}. Second, the literature does not give a
unique sign prediction for the AI-exposure interaction. A
constraints-amplify view would predict larger responses among firms
that are young, leveraged, or rapidly investing --- a description that
fits part of the AI-exposed group but not all of it --- while an
\citet{ottonelloWinberry2020}-style view would predict the opposite in
the leverage dimension. The mechanism narrative through which AI
exposure is expected to matter for equity responses is therefore not
financial heterogeneity, and the next section develops the alternative
asset-pricing channel through growth options, intangible capital, and
long-duration equity.

\section{AI exposure, intangibles, growth options, and long-duration equity}
\label{sec:literature-ai-intangibles}

The closing paragraph of \cref{sec:literature-firm-heterogeneity}
argued that the mechanism through which AI exposure is hypothesized
to matter for firm-level equity responses to forward-guidance shocks
is not financial heterogeneity. This section develops the alternative.
The argument runs in three steps. First, AI engagement is measurable
at the firm level from corporate disclosures, and the resulting
measures are priced by markets in ways consistent with substantive
content rather than pure narrative. Second, AI-investing firms display
the empirical features of growth firms: higher growth, more innovation,
larger market valuations, and elevated systematic risk. Third, growth
firms with these features --- long-duration cash flows, large intangible
capital, and prominent growth options --- are the firms that
asset-pricing theory identifies as disproportionately sensitive to
news about expected discount rates. The chain connects ``AI exposure''
to ``differential response to forward-guidance surprises'' through
established asset-pricing mechanisms rather than through a direct
AI-monetary-policy channel, which the existing literature has not
documented.

\subsection*{Measuring AI engagement from corporate disclosures}

\Citet{anteSaggu2025} provide the most direct methodological precedent
for the AI-exposure measure constructed in this thesis. Using natural
language processing on SEC 10-K filings for 3{,}395 NASDAQ-listed
firms over 2010 to 2022, they construct firm-level AI engagement
metrics from the frequency of AI-related terms in annual disclosures,
adjusted for document length, and validate the resulting classification
by an event study around the November 2022 launch of ChatGPT. Their
central finding is that firms classified as AI-related on the basis
of their 10-K text exhibit significantly larger positive cumulative
abnormal returns following the ChatGPT launch than firms that are not
so classified. The validation is not a test of monetary transmission,
but it establishes that 10-K-based AI engagement measures track
market-perceived AI involvement in a manner that is detectable in
returns. The thesis adopts the same methodological family --- keyword
counts on pre-event 10-K filings, normalized by document length ---
while differing in three respects: the universe is the S\&P 500 rather
than NASDAQ-listed firms; the keyword set is broader, covering
machine learning, deep learning, neural networks, natural language
processing, computer vision, large language models, and generative
AI alongside ``artificial intelligence''; and the exposure measure is
fixed pre-ChatGPT to prevent the AI-boom narrative from contaminating
the classification used to study events that include the post-2022
period. The use of \citet{anteSaggu2025} here is methodological:
their results validate the broader idea that 10-K-based AI exposure is
informative about market-perceived AI involvement, not a claim that
their specific measure equals AI \emph{adoption}.

\Citet{anteSaggu2025} are explicit about the disclosure-vs-adoption
gap. They observe that ``AI'' appears in 10-K filings in heterogeneous
contexts --- some firms describe genuine integration, others mention
competitors' use of AI, and some only state future intentions --- and
they note that the SEC has charged investment advisers for misleading
AI-related claims (``AI washing''). For this thesis, the implication
is that the AI exposure measure should be interpreted as a textual
measure of disclosed AI engagement, capturing some combination of
substantive technological involvement, disclosure choices, sectoral
patterns, and narrative strategy. The pre-2023 timing of the SEC
filings used to build the exposure measure mitigates one specific
concern --- post-ChatGPT strategic AI disclosure --- but it does not
resolve the more general adoption-vs-disclosure distinction. This
caveat is carried forward to the interpretation of the heterogeneity
results in \cref{ch:results}.

\subsection*{AI investment, firm growth, and systematic risk}

A separate literature documents the firm-level fundamentals of AI
investment. \Citet{babinaEtAl2024} construct a firm-level measure of
AI investment from job-posting data on AI-related skill demand and
estimate the effect of AI investment on firm growth using exposure to
universities' supply of AI graduates as an instrument. They find that
AI-investing firms experience higher sales growth, higher employment
growth, and higher market valuations, with product innovation as the
proximate channel. The result establishes AI investment as
empirically associated with the textbook features of growth firms ---
rapid expansion and high valuation multiples --- but it does not by
itself say anything about monetary-policy transmission.

\Citet{babinaEtAlSystematicRisk} extend the analysis to the risk
characteristics of AI-investing firms and find that AI investment
is associated with elevated market beta, measured against the
standard CAPM and confirmed in the Carhart four-factor model.
They interpret this pattern through a growth-options channel and
support it with three concrete pieces of evidence. First,
AI-investing firms experience declining loadings on the
high-minus-low value factor and rising correlation with growth
firms, becoming empirically more growth-like in the value--growth
cross-section. Second, the increase in market beta concentrates
in upside beta --- comovement on days when the market return is
above average --- rather than downside beta, consistent with the
option-like payoff structure of AI-driven product innovation.
Third, the increase in beta is larger on days with elevated AI
news and search attention. Together these results support the
proposition that AI-investing firms identified by 10-K text
measures are not only larger and faster-growing but also more
exposed to systematic risk along the same growth-firm dimension
that the asset-pricing literature in
\cref{sec:literature-equity-valuation} associates with stronger
sensitivity to discount-rate news.

\Citet{eisfeldtSchubertZhangTaska2023} provide a final piece of
validation that AI is priced as a firm characteristic. Using a
workforce-exposure measure constructed from occupational tasks
amenable to generative AI, they document that firms with greater
exposure to generative AI experienced higher returns following the
launch of ChatGPT, with the effect strongest where firms also possess
large data assets. The result is methodologically independent of
\citet{anteSaggu2025} --- different exposure measure, different sample,
different identification --- and arrives at the same qualitative
conclusion: AI exposure, however measured, is a priced firm
characteristic. The \citet{eisfeldtSchubertZhangTaska2023} sample
includes the ChatGPT event itself, which is by construction excluded
from the present thesis's pre-ChatGPT exposure measure; the two
designs are therefore complementary rather than overlapping.

\subsection*{Intangibles, growth options, and long-duration equity}

The asset-pricing bridge from the previous subsection to the
mechanism narrative of this thesis runs through the intangibles and
growth-options literature. \Citet{petersTaylor2017} show that
better measurement of intangible capital substantially strengthens
the empirical investment-Q relation, implying that intangible-intensive
firms are mismeasured by standard balance-sheet metrics in ways that
matter for asset-pricing tests. \Citet{eisfeldtPapanikolaou2013}
provide complementary evidence that firms with more organization
capital --- a form of intangible capital --- have distinct
cross-sectional risk and return characteristics, again implying
that intangible-intensive firms are not interchangeable with
their tangible-capital counterparts in asset-pricing models. AI
investment is plausibly correlated with intangible capital
\citep{petersTaylor2017,eisfeldtPapanikolaou2013}, although the
mapping from textual AI exposure to a balance-sheet intangibles
measure is not one-to-one and is not estimated here.

The asset-pricing implication of intangibles-intensity and
growth-options-intensity is the duration of expected cash flows.
\Citet{campbellPolkVuolteenaho2010} document that growth stocks
carry disproportionately high covariances of returns with market
discount-rate news, while value stocks carry disproportionately
high covariances with market cash-flow news, and they show that
these differential exposures are driven by the cash-flow
fundamentals of growth and value firms rather than by transitory
shifts in investor sentiment. \Citet{vanBinsbergen2020} formalises
the closely related duration-based view of equity valuation: the
sensitivity of a stock's price to a change in the expected
discount-rate path is increasing in the average maturity of its
expected cash flows, in direct analogy with bond duration. The
\citet{campbellPolkVuolteenaho2010} finding is the empirical
counterpart, and the fundamentals-not-sentiment qualification
matters for the mechanism narrative because it locates the source
of growth firms' discount-rate sensitivity in the structure of
their cash flows rather than in investor sentiment about them. 
Putting the chain together: AI-investing firms display
growth-firm characteristics \citep{babinaEtAl2024,babinaEtAlSystematicRisk};
growth firms have longer-duration cash flows
\citep{campbellPolkVuolteenaho2010,vanBinsbergen2020}; and
longer-duration cash flows make a firm's equity price more sensitive
to news about the expected future policy-rate path
\citep{vanBinsbergen2020}. The forward-guidance/path shock used in
this thesis (\cref{sec:literature-forward-guidance}) is the empirical
counterpart to news about the expected discount-rate path, and the
chain therefore predicts that AI-exposed firms should respond more
strongly than non-AI firms to path shocks of a given magnitude, in
the direction consistent with the discount-rate channel.

\subsection*{Status of the mechanism narrative}

Two qualifications are essential. First, the mechanism narrative
above is a synthesis assembled from four distinct literatures rather
than a single causal chain documented in one paper. No paper in the
verified evidence base of this thesis establishes ``AI exposure $\rightarrow$
longer cash-flow duration $\rightarrow$ stronger response to
forward-guidance shocks'' end-to-end. Each link is supported by
existing work, but the composite claim is the thesis's own. Second,
the asset-pricing chain identifies one channel through which AI
exposure could matter; the present empirical design cannot separately
identify it from the other channels in the \citet{campbellAmmer1993}
decomposition. A finding that AI-exposed firms respond more strongly
than non-AI firms to path shocks within sector and event would be
consistent with the discount-rate mechanism described here, but it
would also be consistent with cash-flow-news or risk-premium channels
operating differentially across the AI groups, and the thesis does
not attempt to discriminate among these alternatives.

\section{Gap and contribution}
\label{sec:literature-gap}

The literatures reviewed in
\cref{sec:literature-hfi,sec:literature-forward-guidance,sec:literature-equity-valuation,sec:literature-firm-heterogeneity,sec:literature-ai-intangibles}
form four largely separate strands. The first --- high-frequency
identification of monetary-policy surprises --- has converged on a
small set of widely used shock series, including the forward-guidance
factor used in this thesis \citep{gurkaynakSackSwanson2005,swanson2021}.
The second --- the asset-pricing literature on monetary policy and
equity valuation --- has established that monetary surprises move
equity prices through a composite of discount-rate, cash-flow, and
risk-premium channels, with growth firms disproportionately exposed
to the discount-rate component \citep{bernankeKuttner2005,campbellAmmer1993,campbellPolkVuolteenaho2010}.
The third --- firm-level heterogeneity in monetary transmission ---
has documented that responses differ systematically along financial
and structural characteristics such as leverage, debt maturity, age,
and collateral \citep{cloyneEtAl2023,ottonelloWinberry2020,durante2022,arinEtAl2025}.
The fourth --- AI investment and firm fundamentals --- has shown that
AI-exposed firms are priced as a distinct asset-pricing characteristic
and display the empirical features of growth firms
\citep{anteSaggu2025,babinaEtAl2024,eisfeldtSchubertZhangTaska2023}.

The intersection of these four strands has not been studied. To the
best of the verified evidence base of this thesis, no existing work
estimates firm-level equity responses to forward-guidance shocks
along an AI-exposure heterogeneity dimension, using pre-ChatGPT 10-K
filings to measure exposure and panel local projections to recover
the dynamic response across horizons. Each strand individually
motivates the design --- the shock series is established, the
asset-pricing channels are established, the firm-level heterogeneity
methodology is established, and the AI-exposure measurement methodology
is established --- but the intersection is empirically open.

The contribution of this thesis is therefore primarily empirical: it
adds AI exposure as a new dimension of firm-level heterogeneity in
the literature on monetary transmission to equity prices. Three design
features delimit the scope of this contribution.

First, AI exposure is measured from pre-2023 SEC 10-K filings. This
timing choice prevents the post-ChatGPT AI narrative from contaminating
the classification used to study events that include the post-2022
period. The cost of this choice is that the AI exposure measure is
fixed within firm and does not capture the substantial rise in AI
disclosure documented by \citet{anteSaggu2025} from 2016 onward; it is
a cross-sectional rather than time-varying characteristic. The
benefit is that the classification is predetermined relative to every
FOMC event in the post-ChatGPT subsample, which protects against
look-ahead bias in the heterogeneity estimates.

Second, the empirical design estimates differential responses by AI
exposure group relative to a No AI Exposure reference group, within
sector and FOMC event. The preferred specification absorbs firm
fixed effects and sector-event fixed effects, so that the coefficients
on the AI-exposure interactions measure within-sector, within-event
differential responses rather than unconditional level effects. This
design is conservative: it removes any sector-level response to a
given FOMC event, which is precisely the variation a less demanding
specification would attribute to AI exposure if AI-intensive sectors
respond systematically more to forward-guidance surprises than other
sectors. The implication is that any estimated AI heterogeneity must
be present \emph{within} sector to register in the preferred
specification, which is a higher bar than a cross-sectoral comparison
would impose.

Third, the design does not separately identify the channel through
which AI exposure matters. As emphasized in
\cref{sec:literature-equity-valuation} and again in
\cref{sec:literature-ai-intangibles}, an AI-exposure differential in
the response to forward-guidance surprises could reflect discount-rate-channel
heterogeneity, cash-flow-news heterogeneity, or risk-premium heterogeneity
across the AI groups. The mechanism narrative developed in
\cref{sec:literature-ai-intangibles} is one consistent interpretation
of such a finding --- AI-exposed firms as longer-duration, more
discount-rate-sensitive equity --- but the data do not discriminate
among the channels in the \citet{campbellAmmer1993} decomposition,
and the thesis does not claim to do so.

Three additional features further delimit the contribution. The AI
exposure measure is textual, capturing disclosed AI engagement rather
than AI adoption directly \citep{anteSaggu2025}. The forward-guidance
surprise inherits the information-effect concerns reviewed in
\cref{sec:literature-forward-guidance}; the thesis addresses these
through a path-information split rather than resolving them. The
medium-horizon estimation window (cumulative log returns up to sixty
trading days) extends beyond the announcement-window horizon that is
the focus of most firm-level event studies in this literature; this
extension produces estimates of how long any AI differential persists,
but the longer horizons mix the announcement effect with subsequent
information arrivals and are therefore most cleanly read as
reduced-form persistence rather than causal long-run effects of the
announcement.

The remainder of the thesis is organized accordingly.
\Cref{ch:data} describes the construction of the firm-level panel,
the AI exposure measure, and the shock series.
\Cref{ch:empirical-strategy} develops the panel local-projection design,
states the preferred specification, and discusses identification.
\Cref{ch:results} reports the main impulse-response estimates for the
preferred categorical AI specification, and \cref{ch:robustness}
reports the path-information split and the financial-control
robustness checks discussed in
\cref{sec:literature-forward-guidance,sec:literature-firm-heterogeneity}.
\Cref{ch:conclusion} concludes.